\documentclass[10pt,a4paper]{amsart}
\usepackage{amsmath,amssymb,amsthm,graphicx}
% Retained preamble for the Stacks Project chapter snapshot in this directory.
% The chapter's own preamble.tex is a copy of the shared Stacks Project preamble; the block
% below is sliced MECHANICALLY, line for line, from it. Nothing was added, removed or
% redefined; the wrapper class, the xr-hyper/hyperref wiring and the externaldocument list
% are omitted because the assembler supplies the document class.
%
% lines taken from preamble.tex: 3, 5, 9, 11, 12, 156, 157, 158, ... (51 lines total)
%
% package coverage restored from the chapter's own preamble (2026-10-13): xy, xy-2cell

\usepackage[all]{xy}
\xyoption{2cell}
\UseAllTwocells

\newenvironment{reference}{\comment}{\endcomment}
\newenvironment{slogan}{\comment}{\endcomment}
\newenvironment{history}{\comment}{\endcomment}
\theoremstyle{plain}
\newtheorem{theorem}[subsection]{Theorem}
\newtheorem{proposition}[subsection]{Proposition}
\newtheorem{lemma}[subsection]{Lemma}
\theoremstyle{definition}
\newtheorem{definition}[subsection]{Definition}
\newtheorem{example}[subsection]{Example}
\newtheorem{exercise}[subsection]{Exercise}
\newtheorem{situation}[subsection]{Situation}
\theoremstyle{remark}
\newtheorem{remark}[subsection]{Remark}
\newtheorem{remarks}[subsection]{Remarks}
\numberwithin{equation}{subsection}
\def\lim{\mathop{\mathrm{lim}}\nolimits}
\def\colim{\mathop{\mathrm{colim}}\nolimits}
\def\Spec{\mathop{\mathrm{Spec}}}
\def\Hom{\mathop{\mathrm{Hom}}\nolimits}
\def\Ext{\mathop{\mathrm{Ext}}\nolimits}
\def\SheafHom{\mathop{\mathcal{H}\!\mathit{om}}\nolimits}
\def\SheafExt{\mathop{\mathcal{E}\!\mathit{xt}}\nolimits}
\def\Sch{\mathit{Sch}}
\def\Mor{\mathop{\mathrm{Mor}}\nolimits}
\def\Ob{\mathop{\mathrm{Ob}}\nolimits}
\def\Sh{\mathop{\mathit{Sh}}\nolimits}
\def\NL{\mathop{N\!L}\nolimits}
\def\CH{\mathop{\mathrm{CH}}\nolimits}
\def\proetale{{pro\text{-}\acute{e}tale}}
\def\etale{{\acute{e}tale}}
\def\QCoh{\mathit{QCoh}}
\def\Ker{\mathop{\mathrm{Ker}}}
\def\Im{\mathop{\mathrm{Im}}}
\def\Coker{\mathop{\mathrm{Coker}}}
\def\Coim{\mathop{\mathrm{Coim}}}
\DeclareMathSymbol{\boxtimes}{\mathbin}{AMSa}{"02}
\def\QCohstack{\mathcal{QC}\!\mathit{oh}}
\def\Cohstack{\mathcal{C}\!\mathit{oh}}
\def\Spacesstack{\mathcal{S}\!\mathit{paces}}
\def\Quotfunctor{\mathrm{Quot}}
\def\Hilbfunctor{\mathrm{Hilb}}
\def\Curvesstack{\mathcal{C}\!\mathit{urves}}
\def\Polarizedstack{\mathcal{P}\!\mathit{olarized}}
\def\Complexesstack{\mathcal{C}\!\mathit{omplexes}}
\def\Pic{\mathop{\mathrm{Pic}}\nolimits}
\def\Picardstack{\mathcal{P}\!\mathit{ic}}
\def\Picardfunctor{\mathrm{Pic}}
\def\Deformationcategory{\mathcal{D}\!\mathit{ef}}
\usepackage{cleveref}
\crefname{originallemma}{original lemma}{original lemmas}
\Crefname{originallemma}{Original lemma}{Original lemmas}
\makeatletter
\newlabel{fields-lemma-separable-trace-pairing}{{0BIL}{source}{}{sourcepointer-1}{}}
\newlabel{fields-lemma-separable-trace-pairing@cref}{{[originallemma][0][]0BIL}{[1][0][]source}{}{}{}}
\makeatother
\crefname{originallemma}{original lemma}{original lemmas}
\Crefname{originallemma}{Original lemma}{Original lemmas}
\makeatletter
\newlabel{lemma-base-change-different}{{0BW7}{source}{}{sourcepointer-2}{}}
\newlabel{lemma-base-change-different@cref}{{[originallemma][0][]0BW7}{[1][0][]source}{}{}{}}
\makeatother
\crefname{originallemma}{original lemma}{original lemmas}
\Crefname{originallemma}{Original lemma}{Original lemmas}
\makeatletter
\newlabel{lemma-discriminant}{{0BJF}{source}{}{sourcepointer-3}{}}
\newlabel{lemma-discriminant@cref}{{[originallemma][0][]0BJF}{[1][0][]source}{}{}{}}
\makeatother
\crefname{originallemma}{original lemma}{original lemmas}
\Crefname{originallemma}{Original lemma}{Original lemmas}
\makeatletter
\newlabel{lemma-norm-different-in-discriminant}{{0BW8}{source}{}{sourcepointer-4}{}}
\newlabel{lemma-norm-different-in-discriminant@cref}{{[originallemma][0][]0BW8}{[1][0][]source}{}{}{}}
\makeatother
\crefname{originallemma}{original lemma}{original lemmas}
\Crefname{originallemma}{Original lemma}{Original lemmas}
\makeatletter
\newlabel{morphisms-lemma-set-points-where-fibres-etale}{{0476}{source}{}{sourcepointer-5}{}}
\newlabel{morphisms-lemma-set-points-where-fibres-etale@cref}{{[originallemma][0][]0476}{[1][0][]source}{}{}{}}
\makeatother
\crefname{originallemma}{original lemma}{original lemmas}
\Crefname{originallemma}{Original lemma}{Original lemmas}
\makeatletter
\newlabel{morphisms-section-etale}{{02GH}{source}{}{sourcepointer-6}{}}
\newlabel{morphisms-section-etale@cref}{{[originallemma][0][]02GH}{[1][0][]source}{}{}{}}
\makeatother
\title{Stacks, Discriminants: lemma-different-ramification}
\author{The Stacks Project Authors}
\date{stacks-project:discriminant.tex:local-snapshot-20261012; GNU FDL 1.2 or later}
\begin{document}
\maketitle
\noindent\textit{Editorial note.} Statement span is lines 1906-1912 of sources/stacks-discriminant/pinned-original-source.tex; proof span is lines 1914-1939, the complete author proof (24 body lines).\par\medskip
\noindent\textit{Editorial note.} Navigation stubs: 6 cited-but-not-included labels are supplied as non-load-bearing pointers with their pinned tags where the snapshot tag index has them.\par\medskip
\noindent\textit{Editorial note.} extra\_packages is empty on purpose: the retained preamble already loads xy with options and a second declaration clashes in the assembler wrapper.\par\medskip
\noindent\textit{Editorial note.} Text is retained verbatim, including the author's own typos and placeholders.\par\medskip
\noindent\textit{Editorial note.} Queue index 18 of queue\_discriminant.json (pb=1351, stmt 1906-1912, proof 1914-1939).\par\medskip
\noindent\textit{Allowed standard prerequisites (editorial).}\begin{itemize}
\item Explicit prerequisite boundary. The cited inputs are prerequisites; the argument in the retained span is not replaced by any cited result.
\end{itemize}\medskip
\ifdefined\hypertarget\hypertarget{sourcepointer-1}{}\fi
\noindent\textit{Source reference.} Original lemma 0BIL: sources/stacks-discriminant/pinned-original-source.tex, source lines 1939--1939. Cited by the retained statement or proof but not part of the collected spans; supplied only as a navigation pointer. The argument does not depend on its content.\par\medskip
\ifdefined\hypertarget\hypertarget{sourcepointer-2}{}\fi
\noindent\textit{Source reference.} Original lemma 0BW7: sources/stacks-discriminant/pinned-original-source.tex, source lines 1813--1827. Cited by the retained statement or proof but not part of the collected spans; supplied only as a navigation pointer. The argument does not depend on its content.\par\medskip
\ifdefined\hypertarget\hypertarget{sourcepointer-3}{}\fi
\noindent\textit{Source reference.} Original lemma 0BJF: sources/stacks-discriminant/pinned-original-source.tex, source lines 756--759. Cited by the retained statement or proof but not part of the collected spans; supplied only as a navigation pointer. The argument does not depend on its content.\par\medskip
\ifdefined\hypertarget\hypertarget{sourcepointer-4}{}\fi
\noindent\textit{Source reference.} Original lemma 0BW8: sources/stacks-discriminant/pinned-original-source.tex, source lines 1863--1867. Cited by the retained statement or proof but not part of the collected spans; supplied only as a navigation pointer. The argument does not depend on its content.\par\medskip
\ifdefined\hypertarget\hypertarget{sourcepointer-5}{}\fi
\noindent\textit{Source reference.} Original lemma 0476: sources/stacks-discriminant/pinned-original-source.tex, source lines 1939--1939. Cited by the retained statement or proof but not part of the collected spans; supplied only as a navigation pointer. The argument does not depend on its content.\par\medskip
\ifdefined\hypertarget\hypertarget{sourcepointer-6}{}\fi
\noindent\textit{Source reference.} Original lemma 02GH: sources/stacks-discriminant/pinned-original-source.tex, source lines 1939--1939. Cited by the retained statement or proof but not part of the collected spans; supplied only as a navigation pointer. The argument does not depend on its content.\par\medskip
\begin{lemma}
\label{lemma-different-ramification}
Let $f : Y \to X$ be a flat quasi-finite morphism of Noetherian schemes.
The closed subscheme $R \subset Y$ defined by the different $\mathfrak{D}_f$
is exactly the set of points where $f$ is not \'etale
(equivalently not unramified).
\end{lemma}

\begin{proof}
Since $f$ is of finite presentation and flat, we see that it is \'etale
at a point if and only if it is unramified at that point. Moreover, the
formation of the locus of ramified points commutes with base change.
See Morphisms, Section \ref{morphisms-section-etale} and especially
Morphisms, Lemma \ref{morphisms-lemma-set-points-where-fibres-etale}.
By Lemma \ref{lemma-base-change-different} the formation of $R$ commutes
set theoretically with base change. Hence it suffices to prove the
lemma when $X$ is the spectrum of a field. On the other hand, the
construction of $(\omega_{Y/X}, \tau_{Y/X})$ is local on $Y$.
Since $Y$ is a finite discrete space (being quasi-finite
over a field), we may assume $Y$ has a unique point.

\medskip\noindent
Say $X = \Spec(k)$ and $Y = \Spec(B)$ where $k$ is a field and $B$ is
a finite local $k$-algebra. If $Y \to X$ is \'etale, then
$B$ is a finite separable extension of $k$, and the trace
element $\text{Trace}_{B/k}$ is a basis element of $\omega_{B/k}$
by Fields, Lemma \ref{fields-lemma-separable-trace-pairing}.
Thus $\mathfrak{D}_{B/k} = B$ in this case.
Conversely, if $\mathfrak{D}_{B/k} = B$, then we see from
Lemma \ref{lemma-norm-different-in-discriminant}
and the fact that the norm of $1$ equals $1$ that the
discriminant is empty. Hence
$Y \to X$ is \'etale by Lemma \ref{lemma-discriminant}.
\end{proof}

\end{document}
